% ============================================================
%  preamble.tex — 《对偶、Langlands 对应与 Michael Atiyah 的影响》
%  中文译注本 · 数译 WIT-2021-ATIYAH-LANG
%  编译：xelatex × 2（需要 ctex + fandol + tcolorbox + tikz）
% ============================================================

\usepackage[UTF8,fontset=fandol,heading=true]{ctex}
\usepackage{amsmath,amssymb,mathtools,bm}
\usepackage[a4paper,left=2.7cm,right=2.7cm,top=2.95cm,bottom=2.9cm,headheight=22pt,headsep=9pt]{geometry}
\usepackage{graphicx}
\usepackage{xcolor}
\usepackage{enumitem}
\usepackage{booktabs}
\usepackage{array}
\usepackage{longtable}
\usepackage{tabularx}
\usepackage[font=small,labelfont=bf,labelsep=quad]{caption}
\usepackage[most]{tcolorbox}
\usepackage{tikz}
\usetikzlibrary{arrows.meta,calc,positioning,decorations.markings,shapes.geometric}
\usepackage{fancyhdr}
\usepackage[hidelinks,bookmarksnumbered=true]{hyperref}
\usepackage{footmisc}

% ---------- 颜色 ----------
\definecolor{shuyiBlue}{HTML}{1F4E79}
\definecolor{shuyiGray}{HTML}{595959}
\definecolor{shuyiLight}{HTML}{F5F5F5}
\definecolor{noteFrame}{HTML}{D9822B}
\definecolor{cmtFrame}{HTML}{2E6E8E}
\definecolor{srcFrame}{HTML}{8C8C8C}

% ---------- 版心与行距 ----------
\linespread{1.20}
\setlength{\parskip}{0.28em}
\setlength{\parindent}{2em}
\allowdisplaybreaks[3]
\setlength{\emergencystretch}{2em}

\ctexset{
  section/name={第,章},
  section/number={\chinese{section}},
  section/format={\Large\bfseries\centering},
  section/beforeskip={2.2ex plus .3ex minus .2ex},
  section/afterskip={1.8ex plus .2ex},
  subsection/format={\large\bfseries},
  subsubsection/runin=true,
}

\captionsetup{skip=6pt}

% ---------- 页眉页脚 ----------
\pagestyle{fancy}
\fancyhf{}
\fancyhead[L]{\small\color{shuyiGray}对偶、Langlands 对应与 Michael Atiyah 的影响}
\fancyhead[R]{\small\color{shuyiGray}数译 · 2026}
\fancyfoot[C]{\small\thepage}
\renewcommand{\headrulewidth}{0.3pt}
\renewcommand{\headrule}{\hbox to\headwidth{\color{shuyiGray!45}\leaders\hrule height \headrulewidth\hfill}}
\fancypagestyle{plain}{\fancyhf{}\fancyfoot[C]{\small\thepage}\renewcommand{\headrulewidth}{0pt}}

% ---------- 数学记号 ----------
\newcommand{\CC}{\mathbb{C}}
\newcommand{\RR}{\mathbb{R}}
\newcommand{\ZZ}{\mathbb{Z}}
\newcommand{\PP}{\mathbb{P}}
\newcommand{\cO}{\mathcal{O}}
\newcommand{\cL}{\mathcal{L}}
\newcommand{\cE}{\mathcal{E}}
\newcommand{\cF}{\mathcal{F}}
\newcommand{\cG}{\mathcal{G}}
\newcommand{\cH}{\mathcal{H}}
\newcommand{\cM}{\mathcal{M}}
\newcommand{\cN}{\mathcal{N}}
\newcommand{\cT}{\mathcal{T}}
\newcommand{\cD}{\mathcal{D}}
\newcommand{\cB}{\mathcal{B}}
\newcommand{\cA}{\mathcal{A}}
\newcommand{\dd}{\mathrm{d}}
\newcommand{\ee}{\mathrm{e}}
\newcommand{\ii}{\mathrm{i}}
\DeclareMathOperator{\Hom}{Hom}
\DeclareMathOperator{\Tr}{Tr}
\DeclareMathOperator{\End}{End}
\DeclareMathOperator{\Aut}{Aut}
\DeclareMathOperator{\SL}{SL}
\DeclareMathOperator{\SU}{SU}
\DeclareMathOperator{\SO}{SO}
\DeclareMathOperator{\Sp}{Sp}
\DeclareMathOperator{\GL}{GL}
\DeclareMathOperator{\PSU}{PSU}
\DeclareMathOperator{\ad}{ad}
\DeclareMathOperator{\Gr}{Gr}
\DeclareMathOperator{\Bun}{Bun}
\DeclareMathOperator{\Loc}{Loc}
\newcommand{\gvee}{G^{\vee}}
\newcommand{\Tvee}{T^{\vee}}
\newcommand{\half}{\tfrac{1}{2}}

% 表示、对偶、内积
\newcommand{\rep}[1]{\rho_{#1}}
\newcommand{\ip}[2]{\langle #1,#2\rangle}

% ---------- 底本页码锚 ----------
% 每张幻灯片的起点用一个极小的灰色角标标出，便于回查底本
\newcommand{\srcp}[1]{\textsuperscript{\textcolor{shuyiGray!85}{\tiny\,[原\,p.#1]}}}

% ---------- 章首题记 ----------
\newcommand{\chapterquote}[2]{%
  \begin{flushright}\small\color{shuyiGray}%
  \begin{minipage}{0.86\textwidth}\itshape #1\par\vspace{0.3em}%
  \raggedleft\upshape\footnotesize —— #2\end{minipage}\end{flushright}%
  \vspace{0.6em}}

% ---------- 译注 / 批注 ----------
\newcounter{yizhu}[section]
\renewcommand{\theyizhu}{\arabic{section}.\arabic{yizhu}}
\newcounter{pizhu}[section]
\renewcommand{\thepizhu}{\arabic{section}.\arabic{pizhu}}

\newenvironment{yizhu}[1][]{%
  \par\medskip\refstepcounter{yizhu}%
  \begin{tcolorbox}[enhanced,breakable,colback=noteFrame!4,colframe=noteFrame!60!black,
    boxrule=0.4pt,arc=1.2pt,left=7pt,right=7pt,top=5pt,bottom=5pt,
    fontupper=\small]%
  \noindent{\bfseries\color{noteFrame!55!black}译注 \theyizhu}%
  \ifx\relax#1\relax\else{\bfseries\color{noteFrame!55!black}（#1）}\fi
  \quad\ignorespaces
}{\end{tcolorbox}\par\medskip}

\newenvironment{pizhu}[1]{%
  \par\medskip\refstepcounter{pizhu}%
  \begin{tcolorbox}[enhanced,breakable,colback=cmtFrame!4,colframe=cmtFrame!70!black,
    boxrule=0.4pt,arc=1.2pt,left=7pt,right=7pt,top=5pt,bottom=5pt,
    fontupper=\small]%
  \noindent{\bfseries\color{cmtFrame!80!black}批注 \thepizhu\quad #1}\par\vspace{0.35em}%
}{\end{tcolorbox}\par\medskip}

% ---------- 文献摘录框（英文原文 + 中译） ----------
\newenvironment{abstractbox}[1]{%
  \par\medskip
  \begin{tcolorbox}[enhanced,breakable,colback=shuyiLight,colframe=srcFrame!70,
    boxrule=0.35pt,arc=1pt,left=9pt,right=9pt,top=6pt,bottom=6pt]%
  {\small\bfseries\color{shuyiGray}文献摘录：#1}\par\vspace{0.45em}%
}{\end{tcolorbox}\par\medskip}

% 英文原文用稍小的衬线体
\newcommand{\en}[1]{{\small #1}}

% ---------- 列表 ----------
\setlist[itemize]{topsep=2pt,itemsep=1pt,parsep=0pt,leftmargin=2em}
\setlist[enumerate]{topsep=2pt,itemsep=1pt,parsep=0pt,leftmargin=2.4em}

% ---------- 目录 ----------
\renewcommand{\contentsname}{目\quad 录}
\setcounter{tocdepth}{1}

% ---------- 脚注 ----------
\renewcommand{\footnotesize}{\small}

% ---------- 补充记号（第二次修复） ----------
\newcommand{\HH}{\mathbb{H}}
\newcommand{\OO}{\mathbb{O}}
\newcommand{\cK}{\mathcal{K}}
\newcommand{\slashed}[1]{\not{#1}}
\DeclareMathOperator{\Rep}{Rep}
\DeclareMathOperator{\Perv}{Perv}

% ---------- 公式与图按章编号 ----------
\numberwithin{equation}{section}
\renewcommand{\theequation}{\arabic{section}.\arabic{equation}}
\numberwithin{figure}{section}
\renewcommand{\thefigure}{\arabic{section}.\arabic{figure}}
\numberwithin{table}{section}
\renewcommand{\thetable}{\arabic{section}.\arabic{table}}
